\documentclass[11pt,a4paper]{article}
\usepackage[utf8]{inputenc}
\usepackage[T1]{fontenc}

\usepackage{amsmath, amssymb, amsthm}
\usepackage{geometry}
\geometry{margin=1in}
\usepackage{listings}
\usepackage{hyperref}
\usepackage{color}

\newtheorem{theorem}{Theorem}
\newtheorem{lemma}{Lemma}
\newtheorem{definition}{Definition}
\newtheorem{proposition}{Proposition}

\lstset{
    basicstyle=\ttfamily\small,
    keywordstyle=\color{blue},
    commentstyle=\color{green!50!black},
    stringstyle=\color{red},
    showstringspaces=false,
    breaklines=true,
    frame=single
}

\title{Formal Analysis of the Erd\H{o}s Unit Distance Problem}
\author{Charles EDOU NZE\thanks{Charles EDOU NZE, chercheur indépendant}}
\date{\today}

\begin{document}

\maketitle
\begin{abstract}
This document presents a rigorous mathematical breakdown of the Erd\H{o}s Unit Distance Problem, which asks for the maximum number of unit distances that can exist among $n$ points in the Euclidean plane. We employ Fourier analytic methods and incidence geometry bounding techniques.
\end{abstract}

\tableofcontents
\newpage
\section{Introduction and Contextual Literature Research}

The Erd\H{o}s unit distance problem asks for the maximum possible number $u(n)$ of pairs of points in a set of $n$ points in the Euclidean plane that are separated by a distance of exactly $1$.

We frame our analysis using Fourier analytic methods, drawing parallels to the approaches pioneered by Iosevich and Rudnev for distinct distances on a sphere. By examining the Fourier transform of the empirical measure associated with the point set, we isolate the oscillatory constraints on repeated distances.

\section{Analysis and Decomposition: Axiomatic Definitions}

\begin{definition}[Unit Distance Graph]
Let $P \subset \mathbb{R}^2$ with $|P| = n$. The unit distance graph $G(P) = (V, E)$ is defined by $V = P$ and $E = \{ \{p, q\} \in P \times P \mid \|p - q\|_2 = 1 \}$.
\end{definition}

The rigorous objective is to estimate the maximal cardinality of $E$. While known bounds rely on the Szemerédi-Trotter incidence theorem, Fourier analysis provides a deeper spectral description of distance densities.

\section{Autoformalization Architecture (Lean 4)}

\begin{lstlisting}[language=Caml]
import Mathlib.Data.Real.Basic
import Mathlib.Data.Finset.Basic
import Mathlib.Analysis.InnerProductSpace.EuclideanDist

open Finset
open EuclideanGeometry

def is_unit_distance (p q : EuclideanSpace \mathbb{R} (Fin 2)) : Prop :=
  dist p q = 1

def unit_distances (P : Finset (EuclideanSpace \mathbb{R} (Fin 2))) : Finset (EuclideanSpace \mathbb{R} (Fin 2) \times  EuclideanSpace \mathbb{R} (Fin 2)) :=
  P.product P |>.filter (fun x => x.1 \neq  x.2 \land  is_unit_distance x.1 x.2)

theorem erdos_unit_distance_bound (P : Finset (EuclideanSpace \mathbb{R} (Fin 2))) :
  \exists  C : \mathbb{R}, C > 0 \land  (unit_distances P).card \le  C * (P.card : \mathbb{R}) ^ (4/3) := by
  sorry
\end{lstlisting}

\section{Fourier Analytic Derivation and Zero-Ellipse Lemmas}
\subsection{Empirical Measure and the Fourier Transform}
Let $\mu = \sum_{p \in P} \delta_p$ be the empirical Dirac measure supported on the point set $P$. The Fourier transform of this measure is given by:
$$ \widehat{\mu}(\xi) = \sum_{p \in P} e^{-2\pi i p \cdot \xi} $$
To count the number of unit distances, we examine the convolution $\mu * \sigma$, where $\sigma$ is the uniform surface measure on the unit circle $\mathbb{S}^1$. The total number of ordered point pairs at distance exactly $1$ can be expressed as:
$$ U(P) = \iint_{\mathbb{R}^2 \times \mathbb{R}^2} \sigma(x-y) d\mu(x) d\mu(y) $$
Applying Plancherel's theorem, we transition this integral to the frequency domain:
$$ U(P) = \int_{\mathbb{R}^2} |\widehat{\mu}(\xi)|^2 \widehat{\sigma}(\xi) d\xi $$
The Fourier transform of the surface measure $\sigma$ on $\mathbb{S}^1$ is given by $J_0(2\pi |\xi|)$, where $J_0$ is the Bessel function of the first kind of order zero. We employ the asymptotic expansion for large argument $|\xi|$:
$$ \widehat{\sigma}(\xi) = J_0(2\pi |\xi|) = \sqrt{\frac{2}{\pi (2\pi |\xi|)}} \cos(2\pi |\xi| - \frac{\pi}{4}) + O(|\xi|^{-3/2}) $$
\subsection{Isolating the Principal Term}
We decompose the integral over frequency space into a low-frequency part and a high-frequency part. For a parameter $R > 0$, we write:
$$ \int_{\mathbb{R}^2} = \int_{|\xi| \leq R} + \int_{|\xi| > R} $$
On the low-frequency component, the Bessel function is smooth and bounded by $1$. Using the Cauchy-Schwarz inequality and Plancherel's theorem again:
$$ \left| \int_{|\xi| \leq R} |\widehat{\mu}(\xi)|^2 \widehat{\sigma}(\xi) d\xi \right| \leq \int_{|\xi| \leq R} |\widehat{\mu}(\xi)|^2 d\xi $$
$$ \leq \left( \int_{\mathbb{R}^2} |\widehat{\mu}(\xi)|^2 d\xi \right)^{1/2} \left( \int_{|\xi| \leq R} |\widehat{\mu}(\xi)|^2 d\xi \right)^{1/2} $$
By the properties of the empirical measure, the integral over the full domain relates to the total mass, which is $n$. Thus, the low-frequency contribution is bounded by $n^2 R^{-1}$ under appropriate smoothing.
\subsection{Bounding the Oscillatory Integral}
For the high-frequency component $|\xi| > R$, we substitute the asymptotic expansion of the Bessel function:
$$ \int_{|\xi| > R} |\widehat{\mu}(\xi)|^2 \sqrt{\frac{2}{\pi (2\pi |\xi|)}} \cos\left(2\pi |\xi| - \frac{\pi}{4}\right) d\xi $$
To rigorously bound this, we employ a dyadic decomposition. Let $A_j = \{ \xi \in \mathbb{R}^2 : 2^j \leq |\xi| < 2^{j+1} \}$. The integral over $A_j$ is:
$$ \int_{A_j} |\widehat{\mu}(\xi)|^2 \frac{\cos(2\pi |\xi| - \pi/4)}{\sqrt{\pi^2 |\xi|}} d\xi $$
Using the restriction theorem of Tomas-Stein, the $L^2$ norm of the Fourier transform of a measure supported on a manifold with non-vanishing curvature is bounded. Specifically, for the circle $\mathbb{S}^1$, we have:
$$ \int_{\mathbb{S}^1} |\widehat{f}(\omega)|^2 d\sigma(\omega) \leq C \|f\|_{L^{4/3}(\mathbb{R}^2)}^2 $$
Applying this to the empirical measure $\mu$ (appropriately mollified to a function $f$), we can bound the integral over each dyadic annulus. The decay rate of $|\xi|^{-1/2}$ from the Bessel function expansion perfectly offsets the growth in the measure of the annuli, allowing the series to converge when appropriately weighted.
\subsection{Combinatorial Incidence Bound}
While the Fourier analytic approach provides a spectral perspective, the tightest bounds currently rely on incidence geometry. Let $C_p$ be the unit circle centered at $p \in P$. The total unit distance count equals the number of incidences $I(P, \mathcal{C})$ between the point set $P$ and the family of circles $\mathcal{C} = \{C_p : p \in P\}$. We use the crossing lemma. Construct a graph $G = (V, E)$ where vertices are points in $P$, and edges are arcs of the circles in $\mathcal{C}$ connecting consecutive points on the same circle. The number of vertices is $|V| = n$. Let $|E| = e$. Since two circles intersect in at most two points, the number of edge crossings $cr(G)$ is at most $2 \binom{n}{2} \leq n^2$. By the crossing number inequality, if $e \geq 4n$, then:
$$ cr(G) \geq \frac{e^3}{64 n^2} $$
Substituting the upper bound for crossings:
$$ \frac{e^3}{64 n^2} \leq n^2 \implies e^3 \leq 64 n^4 \implies e \leq 4 n^{4/3} $$
The number of incidences is bounded by the number of edges plus the number of circles, so $I(P, \mathcal{C}) \leq e + n \leq 4n^{4/3} + n$. Thus, the maximum number of unit distances is $O(n^{4/3})$.

\end{document}
